\section{無限等比級数}\label{節:等比級数}
等比数列 $1,r,r^2,\ldots$ の有限和を $S_n=1+r+\cdots+r^{n-1}$ とおく.\,$rS_n=r+r^2+\cdots+r^n$ を引けば,中間の項が消えるため,
\[
 (1-r)S_n=1-r^n,
 \qquad
 S_n=\frac{1-r^n}{1-r}\quad(r\ne1)
\]
を得る.まず有限の項を足した式を確定させ,その後で $n$ を大きくするのが順序である.

$|r|<1$ なら $r^n\to0$ なので,
\[
 \sum_{k=0}^{\infty}r^k
 :=\lim_{n\to\infty}\sum_{k=0}^{n-1}r^k
 =\frac1{1-r}
\]
となる.無限個の項を一度に足しているのではなく,\textbf{部分和の数列の極限}を定義している.一般に $\sum_{k=1}^{\infty}a_k$ が収束するとは,部分和 $\sum_{k=1}^{n}a_k$ が有限の極限値をもつことをいう.

なぜ $|r|<1$ で $r^n\to0$ になるのかも確認できる.\,$0<q=|r|<1$ とし,\,$\frac1q=1+c$,\,$c>0$ とおく.二項定理から $(1+c)^n\ge1+nc$ なので,
\[
 0\le |r^n|=q^n\le\frac1{1+nc}\longrightarrow0.
\]
$r=0$ の場合は直接明らかである.

$r=1$ なら部分和は $n$ となり,\,$r=-1$ なら部分和は $1,0,1,0,\ldots$ と振動する.\,$|r|>1$ では各項 $r^n$ さえ $0$ に近づかないため,級数は収束しない.級数が収束するなら,隣り合う部分和の差である一般項は必ず $0$ に近づく.この条件は必要条件であり,一般の級数では十分条件ではない.たとえば調和級数の発散は\chapterref{章:積分と数列}で積分との比較から調べる.

\begin{example}{無限和と誤差}
$\sum_{k=0}^{\infty}\left(-\frac13\right)^k$ の和を求めよ.また,\,$k=0$ から $n-1$ までの和で近似するとき,誤差の絶対値を $10^{-4}$ 未満にする自然数 $n$ を一つ挙げよ.
\end{example}
公比の絶対値は $\frac13<1$ だから,和は
\[
 S=\cfrac{1}{1+\cfrac13}=\frac34.
\]
有限和と極限値との差は,有限和の公式から
\[
 S-S_n=\cfrac{\left(-\cfrac13\right)^n}{1+\cfrac13},
 \qquad
 |S-S_n|=\frac{3}{4\cdot3^n}.
\]
$n=9$ では $\frac3{4\cdot19683}<10^{-4}$ となるので,\,$9$ 項まで足せばよい.交互に符号が変わるため誤差の符号も変わるが,誤差の大きさはこの一つの式で評価できる.


